\section{Mathematical and computational background}

\subsection{Lattices}
The fundamental quantity in lattice-based cryptography is a lattice. For the purposes of this text, the following definition is sufficient. 
\begin{definition}[Lattice]
    A lattice $\LL$ is a discrete additive subgroup of $\RR^\nd$.
\end{definition} 
A lattice is finitely generated by a set of linearly independent basis vectors $\{ \vv b_1, \ldots, \vv b_k \}$, where each $\vv b_i \in \RR^k$, often collected in a basis matrix $\vv{B}$ such that
\[ \LL = \LL(\vv{B}) = \vv{B} \cdot \ZZ^k = \left\{ \sum_{i = 1}^{k} z_i \vv b_i : z_i \in \ZZ\right\}. \]
The number of linearly independent basis vectors $k$ is the \textit{rank} of the lattice. We are only intested in lattices of full rank, where $k = \nd$. A characteristic property of a given lattice $\LL$ are it's minima, which relate to the geometry of the lattice. This means the lattice must be endowed with a norm $|| \cdot ||$. Unless otherwise specified, we work in the Euclidean norm.

\begin{definition}[Successive minima]
    The $i$th minimum $\lambda_i(\LL)$ is the radius of the smallest ball centred in the origin containing $i$ linearly independent lattice vectors:
    \[ \lambda_i(\LL) = \inf\{ r : \dim (\mathrm{span} (\LL \cap \mathcal{B}(0,r))) \geq i)\}, \] where $\mathcal{B}(\vv c,r)$ denotes a ball centred on $\vv c$ with radius $r$.
\end{definition}

In particular, the \textit{first minimum}, also known as the \textit{minimum distance} or simply \textit{minimum} is the length of the shortest nonzero lattice vector 
    \[ \lambda_1(\LL) = \min_{v \in \LL \setminus \{ 0 \} } || \vv v ||. \]
Every lattice $\LL$ induces a dual lattice defined as follows.
\begin{definition}[Dual lattice]
    The dual $\LL^*$ of a lattice $\LL$ is defined as
    \[ \LL^* = \{ \vv u : \langle \vv u, \LL \rangle \subseteq \ZZ \}. \]
    In words, the dual $\LL^*$ is the set of points having integer inner product with vectors in $\LL$. 
\end{definition}
The dual $\LL^*$ is also a lattice. If $\vv{B}$ is a basis $\LL$, when $(\vv{B}^T)^{-1}$ is a basis of $\LL^*$ ($T$ denotes transpose).

\subsection{Computational problems on lattices}
We cover the computational problems that have proven useful for constructing cryptography. This is not an exhaustive list of available computational problems. Note that all problems below are approximate problems involving an approximation/gap parameter $\gamma(\nd)$. Problem hardness decreases with $\gamma$, while exact instances of the problems are given by setting $\gamma = 1$.

\begin{definition}[Approximate Shortest Vector Problem ($\gamma$-SVP)]
    Given a lattice basis $\vv{B}$ of an $\nd$-dimensional lattice $\LL(\vv{B})$, find a nonzero lattice vector $v \in \LL$ such that $||\vv v|| \leq \gamma(\nd) \lambda_1(\LL)$.
\end{definition}
In particular, setting $\gamma = 1$ gives the exact Shortest Vector Problem (SVP) of finding a shortest nonzero lattice vector. Note the usage of \textit{a} shortest vector here, there may exist more than one vector $v \in \LL$ that satisfies the condition. SVP is intuitive and likely the most studied computational problem in lattice theory. However, as of yet, no cryptosystem has been proven secure assuming hardness of $\gamma$-SVP. Cryptographic proofs instead reduce to the problems given below.

\begin{definition}[Decisional Approximate SVP ($\gamma$-GapSVP)]
    Given a basis $\vv{B}$ of an $\nd$-dimensional lattice $\LL(\vv{B})$, where either $\lambda_1(\LL) \leq 1$ or $\lambda_1(\LL) > \gamma(n)$, determine which is the case.
\end{definition}

As the name suggests, this is a decision problem (also referred to as a promise problem), in comparison to $\gamma$-SVP, which is a search problem. A slight generalisation of $\gamma$-SVP is the below search problem.

\begin{definition}[Approximate Shortest Independent Vectors Problem ($\gamma$-SIVP)]
    Given a full-rank basis $\vv{B}$ of an $\nd$-dimensional lattice $\LL(\vv{B})$, output a set $S = \{ \vv s_i \} \subset \LL$ of $\nd$ linearly independent lattice vectors where $|| \vv s_i || \leq \gamma(n) \lambda_n(\LL) \; \forall i$.
\end{definition}

Another promise problem is the following, in which we are promised that there exists a lattice point fairly close to a target point $\vv t$. 

\begin{definition}[Bounded Distance Decoding ($\gamma$-BDD)]
    Given a basis $\vv{B}$ of an $\nd$-dimensional lattice $\LL(\vv{B})$ and a target point $\vv t \in \RR^\nd$ such that $\mathrm{dist}(\vv t, \LL) < d = \lambda_1(\LL)/(2 \gamma(\nd))$, find the unique lattice vector $\vv v \in \LL$ closest to the target point $t$, i.e. such that $|| \vv t - \vv v || < d$.
\end{definition}

$\gamma$-BDD is actually a relaxation of the $\gamma$-Closest Vector Problem ($\gamma$-CVP), where the promised existence of a lattice point is removed. CVP is perhaps the second most studied computational lattice problem with a plethora of hardness results. Still however, there is no cryptographic assumption that reduces to CVP at the time of writing. It is intuitively clear that $\gamma$-BDD is easier than $\gamma$-CVP.

The above problems are normally hard to solve given a random basis but significantly easier to solve given a basis of short and close to orthogonal basis vectors. The computational problems above are therefore usually solved using \textit{reduction algorithms}, which turns a ``bad'' basis (containing long and non-orthogonal basis vectors) into a good one. Despite extensive study for several decades, the problems appear to be intractable. There exist polynomial-time algorithms for solving instances with approximation factors $\gamma = 2^{\Theta(\nd \log \log \nd / \log \nd)}$ that are slightly sub-exponential, such as LLL (due to Lenstra, Lenstra and Lovász) and derivatives. For smaller (i.e. polynomial) approximation factors, the algorithms have at least exponential time and/or space complexity.

Some lattice problems are known to be NP-hard 
\subsection{Discrete Gaussians}
Definition.